\documentclass[11pt,a4paper]{article}

\usepackage{lmodern}
\usepackage[T1]{fontenc}
\usepackage[utf8]{inputenc}
\usepackage[margin=2.6cm]{geometry}
\usepackage{amsmath}
\usepackage{amssymb}
\usepackage{booktabs}
\usepackage{microtype}
\usepackage{algorithm}
\usepackage{algpseudocode}
\usepackage[round]{natbib}
\usepackage[hidelinks]{hyperref}

\newcommand{\code}[1]{\texttt{#1}}
\newcommand{\R}{\mathbb{R}}
\DeclareMathOperator{\sign}{sign}
\DeclareMathOperator{\med}{med}

\title{\code{TinyCTA}: An Analytically Scaled Trend Signal and
Correlation-Aware Position Sizing in Nine Modules}

\author{Thomas Schmelzer\thanks{\url{https://github.com/tschm/TinyCTA}}}

\date{\today}

\begin{document}

\maketitle

\begin{abstract}
\code{TinyCTA} is a deliberately small library for building trend-following
strategies: roughly 1{,}300 lines covering signal generation, robust volatility
estimation, and a position-sizing engine that solves a shrunk correlation system
at every timestamp. This paper documents the pipeline and derives the one piece
of it that is stated in the code without proof --- the analytic scale factor of
the oscillator. We show it is exactly the standard deviation of the EWMA
difference under a driftless random walk with unit-variance increments, because
the difference filter's coefficients sum to zero and its cumulative
representation telescopes; the derivation takes four lines, and a simulation over
$2 \times 10^5$ steps confirms unit output standard deviation across a decade of
$(\mathrm{fast}, \mathrm{slow})$ pairs. We then run the full engine over a
$49$-asset, $53$-year futures panel and report what the shrinkage parameter
actually does. It is a numerical necessity rather than a performance lever: at
$\lambda = 1$, where the correlation matrix is used unmodified, the engine raises
\code{SingularMatrixError} --- the empirical correlation matrix over $47$ assets
has condition number $2.8 \times 10^{18}$ --- and as $\lambda \to 1$ mean gross
exposure grows elevenfold while the Sharpe ratio stays inside $[0.21, 0.33]$.
We also flag that the parameter reads backwards: \code{shrink=1.0} means
\emph{no} shrinkage, and it is the setting that fails.
\end{abstract}

\section{Introduction}

A trend-following CTA is a short pipeline. Prices become a signal; the signal is
treated as an expected return; expected returns become positions, sized by
volatility and by the correlation structure of the universe; positions become
P\&L. Each stage is a few lines of arithmetic, and the difficulty is not in any
one of them but in the accumulated conventions --- which returns, which lookback,
which normalisation, what happens to an asset that has not started trading yet.

\code{TinyCTA} is an attempt to write that pipeline down at minimum size and to
keep every convention explicit. It is nine modules \citep{tinycta}: two signal
generators, volatility and matrix utilities, a validated configuration object, a
Polars-facing engine \citep{polars}, a pure-NumPy numeric kernel
\citep{harris2020}, and an optional Optuna \citep{akiba2019} layer for parameter
search. The core has four runtime dependencies and no plotting, no data loading,
and no backtest framework.

The empirical literature on time-series momentum is large
\citep{moskowitz2012,hurst2017,baltas2013} and this paper adds nothing to it.
The contribution here is narrower and, we think, more useful: a precise account
of what this implementation computes, a derivation of the one formula it asserts,
and a measurement of the one parameter whose effect is not obvious.

\section{Signals}

Both signal generators are Polars expressions, so they compose inside a single
\code{with\_columns} call and run column-wise over a wide price frame.

\subsection{Moving-average crossover}

The discrete signal is the sign of an EWMA difference,
\begin{equation}
  s_t = \sign\bigl( \mathrm{EWMA}_{\mathrm{fast}}(x)_t - \mathrm{EWMA}_{\mathrm{slow}}(x)_t \bigr)
  \in \{-1, 0, +1\},
\end{equation}
with the lengths interpreted as centres of mass, $\mathrm{com} = \mathrm{length}
- 1$, so that the decay factor is $1 - 1/\mathrm{length}$. It is computed with
\code{adjust=False}, the recursive form, and emits nulls until
\code{min\_samples} observations have accumulated.

\subsection{The oscillator and its scale factor}

The continuous signal is the same difference, divided by a constant:
\begin{equation}
  \mathrm{osc}_t
  = \frac{\mathrm{EWMA}_{\mathrm{fast}}(x)_t - \mathrm{EWMA}_{\mathrm{slow}}(x)_t}{s},
  \qquad
  s = \sqrt{\frac{1}{1-f^2} - \frac{2}{1-fg} + \frac{1}{1-g^2}},
  \label{eq:osc}
\end{equation}
where $f = 1 - 1/\mathrm{fast}$ and $g = 1 - 1/\mathrm{slow}$. Here the EWMAs use
\code{adjust=True}, the normalised form, whose steady-state weights are
$(1-f)f^k$.

The code states $s$ without derivation. It is worth having, because it is what
makes oscillator magnitudes comparable across parameter choices --- the property
the whole design rests on.

\paragraph{Claim.}
Let $x$ be a driftless random walk, $x_t = \sum_{j \le t} \varepsilon_j$ with
$\varepsilon_j$ uncorrelated of unit variance. Then the numerator of
Eq.~\eqref{eq:osc} has standard deviation exactly $s$.

\paragraph{Proof.}
In steady state the numerator is a linear filter of the price history,
\begin{equation}
  y_t = \sum_{k \ge 0} \bigl[ (1-f)f^k - (1-g)g^k \bigr] x_{t-k},
\end{equation}
whose coefficients sum to $1 - 1 = 0$. A filter with coefficients summing to zero
annihilates the unit root, so $y_t$ is stationary even though $x_t$ is not.
Rewriting in terms of the increments and collecting the coefficient of
$\varepsilon_{t-m}$ gives a telescoping partial sum,
\begin{equation}
  c_m = \sum_{k=0}^{m} \bigl[ (1-f)f^k - (1-g)g^k \bigr]
      = \bigl(1 - f^{m+1}\bigr) - \bigl(1 - g^{m+1}\bigr)
      = g^{m+1} - f^{m+1},
\end{equation}
so that $y_t = \sum_{m \ge 0} c_m \varepsilon_{t-m}$ and
\begin{equation}
  \operatorname{Var}(y_t)
  = \sum_{m \ge 1} \bigl( g^{m} - f^{m} \bigr)^2
  = \frac{g^2}{1-g^2} - \frac{2fg}{1-fg} + \frac{f^2}{1-f^2} .
\end{equation}
Adding and subtracting the $m=0$ term, $(1 - 2 + 1) = 0$, turns each
$z^2/(1-z^2)$ into $1/(1-z^2)$ and recovers $s^2$ exactly as written in
Eq.~\eqref{eq:osc}. \hfill$\square$

So the oscillator is a $t$-statistic of trend: numerator in units of price
moves, denominator the standard deviation those moves would have if the price
were a random walk. Under the null of no trend its output has unit variance,
whatever \code{fast} and \code{slow} are, which is what allows a signal
parameterised at $(8, 24)$ and one at $(64, 192)$ to be added, compared or
blended without a per-parameter rescaling.

\begin{table}[t]
\centering
\caption{Analytic scale factor against the realised standard deviation of the
oscillator on a simulated driftless random walk with unit-variance Gaussian
increments ($2 \times 10^5$ steps, seed $0$). The analytic factor varies by more
than a factor of five across these parameter pairs; the output does not.}
\label{tab:osc}
\begin{tabular}{rrrr}
\toprule
fast & slow & $s$ (analytic) & $\operatorname{sd}(\mathrm{osc})$ \\
\midrule
$2$  & $6$   & $1.0851$ & $0.9983$ \\
$4$  & $12$  & $1.4651$ & $1.0008$ \\
$8$  & $24$  & $2.0334$ & $1.0053$ \\
$16$ & $48$  & $2.8513$ & $1.0106$ \\
$32$ & $96$  & $4.0159$ & $1.0133$ \\
$64$ & $192$ & $5.6680$ & $1.0083$ \\
$8$  & $96$  & $6.1323$ & $1.0117$ \\
\bottomrule
\end{tabular}
\end{table}

Table~\ref{tab:osc} is the numerical check. The realised standard deviation sits
within $1.4\%$ of unity for every pair tested, including the deliberately
mismatched $(8, 96)$. The small positive bias is a finite-sample effect: the
steady-state weights assumed above are only approached after the slower EWMA has
warmed up.

\section{Volatility and matrix utilities}

\paragraph{Volatility-adjusted returns.}
\code{vol\_adj} standardises log returns by an EWMA standard deviation and clips
the result,
\begin{equation}
  \tilde{r}_t = \operatorname{clip}\!\left( \frac{r_t}{\hat{\sigma}_t},\; -c,\; +c \right),
  \qquad r_t = \log x_t - \log x_{t-1},
\end{equation}
with $\hat\sigma$ an EWMA of $r$ at centre of mass $\mathrm{vola}-1$. The clip
is what keeps a single gap or a bad print from dominating everything downstream;
the default in the engine is $c = 4.2$. \code{adj\_log\_prices} integrates
$\tilde{r}$ by cumulative sum, giving a price-like series of roughly unit
volatility --- the natural input when a signal should see moves in risk units
rather than in currency.

Note the warmup: an EWMA standard deviation is undefined for a single
observation, so the first log return produces no value and the output starts at
the second, irrespective of \code{min\_samples}.

\paragraph{Robust volatility.}
\code{moving\_absolute\_deviation} is a doubly-robust alternative: both the
centre and the dispersion are rolling medians of log returns,
\begin{equation}
  \hat\sigma^{\mathrm{MAD}}_t
  = \frac{1}{0.6745} \,
    \med_{w}\Bigl( \bigl| r_t - \med_{w}(r_t) \bigr| \Bigr),
  \qquad w = 2\,\mathrm{com} - 1 .
\end{equation}
The constant $0.6745$ is the third quartile of the standard normal, so the
estimator is consistent for the standard deviation under normality while
retaining the median's breakdown point of $50\%$ \citep{huber2009}. Chaining two
rolling medians over a differenced series costs $2w - 1$ returns of warmup before
the first value appears.

\paragraph{Shrinkage towards the identity.}
\begin{equation}
  \mathrm{shrink2id}(M, \lambda) = \lambda M + (1-\lambda) I .
  \label{eq:shrink}
\end{equation}
For a correlation matrix this preserves the unit diagonal and pulls every
off-diagonal entry towards zero in proportion to $1-\lambda$; it is the
identity-target special case of the classical shrinkage estimator
\citep{ledoit2004}.

\paragraph{A warning about the parameter's direction.}
Equation~\eqref{eq:shrink} reads $\lambda = 1 \Rightarrow$ no shrinkage and
$\lambda = 0 \Rightarrow$ complete shrinkage, and the engine's configuration
field carrying $\lambda$ is named \code{shrink}. So \code{shrink=1.0} applies
\emph{no} shrinkage and \code{shrink=0.0} discards the correlations entirely ---
the opposite of how the name reads. Section~\ref{sec:shrink} shows this is not a
cosmetic complaint: \code{shrink=1.0} is exactly the setting that fails on real
data.

\section{The position engine}

\subsection{Configuration}

Four parameters, validated by a frozen Pydantic model that rejects unknown keys:
\code{vola} and \code{corr} (the two lookbacks, both positive), \code{clip}
(positive), and \code{shrink} $\in [0,1]$. One cross-field constraint is
enforced: $\mathrm{corr} \ge \mathrm{vola}$, on the grounds that estimating a
correlation matrix over a shorter window than the volatilities it is built from
is numerically unstable. Freezing the model means a validated configuration
cannot drift between construction and use.

\subsection{Per-timestamp sizing}

The engine consumes two aligned wide frames --- \code{prices} and \code{mu}, the
expected returns, with identical shapes and columns and a \code{date} column in
each --- and produces a frame of cash positions of the same shape.

Internally it derives clipped volatility-adjusted returns, an EWMA covariance
matrix per timestamp over a window of $2\,\mathrm{corr} + 1$ with $\mathrm{corr}$
observations of warmup, and normalises each to a correlation matrix $C_t$ by its
own diagonal. A zero-variance asset yields NaN rather than a division by zero.
The forward walk is then delegated to a pure-NumPy kernel.

At timestamp $t$, writing $\mathcal{A}_t$ for the set of assets with a finite
price (so an instrument that has not started trading is absent, not zero),
$\mu_t$ for the expected returns restricted to $\mathcal{A}_t$ with NaNs zeroed,
and $\rho_t = \mathrm{shrink2id}(C_t, \lambda)$ restricted likewise:
\begin{align}
  u_t &= \frac{\rho_t^{-1} \mu_t}{\sqrt{\mu_t^{\top} \rho_t^{-1} \mu_t}},
    \label{eq:risk} \\
  \pi_t &= \frac{u_t}{v_t},
    \label{eq:scale} \\
  h_{t,i} &= \frac{\pi_{t,i}}{\hat\sigma_{t,i}} .
    \label{eq:cash}
\end{align}

Equation~\eqref{eq:risk} is the informative step. It is a Markowitz direction
\citep{markowitz1952} in correlation space, normalised so that
$u_t^{\top} \rho_t u_t = 1$ --- one unit of risk under the metric $\rho_t$,
exactly. The normaliser is computed as
$\sqrt{\mu_t^{\top}\rho_t^{-1}\mu_t}$, and the position is zeroed rather than
scaled when that quantity is non-finite, below $10^{-12}$, or when $\mu_t$
vanishes.

Equation~\eqref{eq:scale} is a feedback loop. Realised profit
$p_t = h_{t-1}^{\top} r_t$ drives an EWMA of squared profit,
\begin{equation}
  v_t = \kappa \, v_{t-1} + (1-\kappa) \, p_t^2,
  \qquad \kappa = 0.99, \quad v_0 = 1,
  \label{eq:pv}
\end{equation}
and dividing by it shrinks the book after volatile P\&L and expands it after
quiet P\&L --- a volatility target expressed in realised profit rather than in
forecast risk. Note that $\kappa$ and $v_0$ are hard-coded, not configuration.

Equation~\eqref{eq:cash} converts risk units to cash by dividing by each asset's
own EWMA volatility of simple returns, so a given risk allocation buys fewer
units of a volatile instrument.

\paragraph{Scale is not pinned down.}
Equation~\eqref{eq:pv} carries units. $p_t$ is a currency amount, so $v_t$ is a
squared currency amount, while $u_t$ in Eq.~\eqref{eq:risk} is dimensionless;
dividing one by the other leaves the overall size of the book determined jointly
by $v_0 = 1$ and the numerical scale of the price series. There is no
target-volatility or AUM parameter. The gross exposures in
Table~\ref{tab:shrink} should therefore be read as relative, not as leverage
against a stated capital base --- a caller who needs a specific risk level must
rescale the output.

\section{Empirical behaviour}
\label{sec:empirical}

We run the pipeline end to end on the futures panel shipped with the test suite:
$13{,}909$ daily observations of $49$ instruments with hashed identifiers,
spanning 1970-01-02 to 2023-04-26, with instruments entering over time (the
median instrument has about $9{,}100$ observations; all $49$ are live at the
end). Expected returns are $\mathrm{osc}(8, 24)$ per asset, nulls filled with
zero. The configuration is $\mathrm{vola}=32$, $\mathrm{corr}=64$,
$\mathrm{clip}=4.2$, and $\lambda$ varies. P\&L is
$p_t = \sum_i h_{t-1,i} r_{t,i}$ with $r$ simple returns; the Sharpe ratio is
$\operatorname{mean}(p)/\operatorname{sd}(p) \times \sqrt{252}$, gross of all
costs.

\subsection{What shrinkage does}
\label{sec:shrink}

\begin{table}[t]
\centering
\caption{The engine over the $49$-asset panel as $\lambda$ (\code{shrink})
varies. Recall $\lambda = 1$ means the correlation matrix is used unmodified.
Gross exposure is $\sum_i |h_{t,i}|$, in the units of the price series.}
\label{tab:shrink}
\begin{tabular}{lrrr}
\toprule
$\lambda$ & Sharpe & Mean gross & Max gross \\
\midrule
$0.00$ & $0.265$ & $128.9$   & $706.3$ \\
$0.25$ & $0.225$ & $428.4$   & $4{,}436.8$ \\
$0.50$ & $0.220$ & $629.7$   & $7{,}423.6$ \\
$0.75$ & $0.255$ & $889.2$   & $9{,}891.2$ \\
$0.90$ & $0.327$ & $1{,}168.7$ & $12{,}229.1$ \\
$0.99$ & $0.206$ & $1{,}415.4$ & $219{,}767.6$ \\
$1.00$ & \multicolumn{3}{c}{\code{SingularMatrixError}} \\
\bottomrule
\end{tabular}
\end{table}

Table~\ref{tab:shrink} contains the paper's main empirical finding, and it is
about numerics rather than returns.

At $\lambda = 1$ the engine does not produce a worse result --- it produces no
result. The solve behind Eq.~\eqref{eq:risk} raises
\code{SingularMatrixError}: the Cholesky factorisation fails as not
positive-definite and the fallback general solve reports a singular matrix. This
is not a fluke of one timestamp. The final EWMA correlation matrix over the $47$
assets with finite entries has condition number $2.79 \times 10^{18}$ and a
smallest eigenvalue of $9.16 \times 10^{-17}$ --- numerically rank-deficient at
double precision, notwithstanding an EWMA window ($129$) considerably longer
than the cross-section ($49$). Clipping and the overlapping decay weights leave
the effective sample far smaller than the nominal window.

Shrinkage is therefore load-bearing infrastructure in this engine, not a tuning
knob: any $\lambda < 1$ makes $\rho_t$ invertible, because
Eq.~\eqref{eq:shrink} adds $(1-\lambda)$ to every eigenvalue.

The route to that failure is visible in the exposure columns, and it is the more
practically useful observation. Mean gross exposure rises elevenfold from
$\lambda = 0$ to $\lambda = 0.99$, and maximum gross exposure rises by a factor
of $311$, reaching $2.2 \times 10^5$ at $\lambda = 0.99$ against $706$ at
$\lambda = 0$. Ill-conditioning does not announce itself as an exception until
the very end; long before that it shows up as leverage, as
$\rho_t^{-1}$ amplifies small differences between nearly collinear assets
\citep{michaud1989}. A practitioner monitoring positions would see it; one
monitoring only the Sharpe ratio would not.

Because the Sharpe ratio does not see it. Across the whole feasible range it
stays inside $[0.206, 0.327]$, non-monotone in $\lambda$, with the best value at
$\lambda = 0.90$ --- and we would caution explicitly against reading that as a
recommendation. The spread is well inside what $53$ years of one gross-of-costs
configuration can distinguish \citep{harvey2018}, and the $\lambda = 0.99$
column shows how a nominally similar Sharpe can be produced by a book two orders
of magnitude larger at its peak. The honest summary is that correlation shrinkage
buys invertibility and exposure control here, and that this naive configuration
--- one signal, no costs, no turnover control, in-sample parameters --- has no
evidence to offer about its effect on returns.

\subsection{Parameter search}

The optional \code{tinycta.hyper} layer wraps Optuna \citep{akiba2019} with a
tree-structured Parzen estimator \citep{bergstra2011}, maximising the Sharpe
ratio of a caller-supplied portfolio-suggestion function. Trials returning NaN
are pruned rather than failed, which keeps the reported statistics clean. The
layer is behind an extra (\code{pip install "tinycta[hyper]"}) and imports
nothing that the core needs; the numeric kernel deliberately avoids even a
logging dependency so that a core install stays at four packages.

Given the width of the Sharpe band in Table~\ref{tab:shrink}, we note that the
same caution applies with more force to any search over more parameters: an
optimiser that maximises in-sample Sharpe over a four-dimensional configuration
will find something, and this panel cannot tell whether it found a strategy.

\section{Limitations}

\paragraph{No trading frictions.} No transaction costs, no slippage, no market
impact, no financing, and no turnover penalty. Positions can change arbitrarily
between consecutive timestamps, and Eq.~\eqref{eq:risk} has no term that
discourages it. Any Sharpe ratio quoted here is gross.

\paragraph{No risk limits.} There are no position caps, no gross or net exposure
limits, no drawdown control, and no leverage constraint. The only mechanism
restraining size is the profit-variance feedback of Eq.~\eqref{eq:pv}, which
reacts to realised P\&L after the fact and, as Table~\ref{tab:shrink} shows,
permits gross exposure to move over two orders of magnitude.

\paragraph{Memory.} The engine materialises one correlation matrix per
post-warmup timestamp in a dictionary. On this panel that is roughly $13{,}800$
matrices of $49 \times 49$ doubles --- about $265$ MB --- and it grows as $O(T
n^2)$. The design is fine at CTA scale and will not survive a large
cross-section or intraday data without being reworked into a streaming walk.

\paragraph{Hard-coded constants.} The profit-variance decay $\kappa = 0.99$, its
initial value $v_0 = 1$, and the degeneracy floor $10^{-12}$ are not
configurable. The first two set the book's scale and its responsiveness, so they
are choices a user might reasonably want to make.

\paragraph{In-sample, single panel.} Every number in
Section~\ref{sec:empirical} comes from one dataset with parameters chosen before
looking at the result but evaluated on the whole history. Nothing here is a
backtest with an out-of-sample period, and nothing here should be read as
evidence that this configuration works.

\section{Conclusion}

The pipeline is small enough to state completely, which is the point of the
package and the point of this paper. Two results are worth carrying away.

The oscillator's scale factor is exact, not heuristic: it is the standard
deviation of the EWMA difference under a random walk, provable in four lines
because the difference filter's coefficients telescope, and confirmed to within
$1.4\%$ in simulation across parameter pairs whose raw magnitudes differ by more
than five times. Signals at different speeds really are comparable without
further normalisation.

Correlation shrinkage is not optional. At $\lambda = 1$ --- which, confusingly,
is what \code{shrink=1.0} means --- the engine cannot run at all on a $49$-asset
panel whose correlation matrix is numerically singular, and as $\lambda$
approaches $1$ the deterioration shows up first as an eleven-fold growth in gross
exposure rather than as an error or a worse Sharpe ratio. The parameter's name
inverts its meaning, and its safe region is bounded away from the value the name
suggests is safest.

\appendix

\section{Reproducibility}
\label{app:repro}

All numbers were produced with \code{TinyCTA} 0.14.0 and
\path{tests/tinycta/resources/prices_hashed.csv}. Columns with long leading gaps
are inferred as strings by default, hence the explicit cast.

\begin{verbatim}
import numpy as np, polars as pl
from tinycta.config import Config
from tinycta.engine import Engine
from tinycta.osc import osc

raw = pl.read_csv("tests/tinycta/resources/prices_hashed.csv",
                  try_parse_dates=True, infer_schema_length=None)
raw = raw.rename({raw.columns[0]: "date"})
assets = [c for c in raw.columns if c != "date"]
prices = raw.with_columns(pl.col(a).cast(pl.Float64) for a in assets)
mu = prices.with_columns(osc(pl.col(a), fast=8, slow=24).alias(a)
                         for a in assets)
mu = mu.with_columns(pl.col(a).fill_null(0.0).fill_nan(0.0) for a in assets)

P = prices.select(assets).to_numpy()
ret = np.zeros_like(P); ret[1:] = P[1:] / P[:-1] - 1.0

for shrink in (0.0, 0.25, 0.5, 0.75, 0.9, 0.99, 1.0):
    cfg = Config(vola=32, corr=64, clip=4.2, shrink=shrink)
    engine = Engine(prices=prices, mu=mu, cfg=cfg)
    H = engine.cash_position.select(assets).to_numpy()
    p = np.nansum(np.nan_to_num(H[:-1]) * np.nan_to_num(ret[1:]), axis=1)
    gross = np.nansum(np.abs(H), axis=1)
    print(shrink, p.mean() / p.std() * np.sqrt(252),
          np.nanmean(gross), np.nanmax(gross))
\end{verbatim}

The $\lambda = 1.0$ row raises \code{cvx.linalg.core.exceptions.\allowbreak
SingularMatrixError} from inside \code{Engine.cash\_position}. Condition numbers
are of the last correlation matrix returned by \code{Engine.cor}, restricted to
rows and columns that are finite throughout. Table~\ref{tab:osc} simulates
\code{np.random.default\_rng(0).standard\_normal(200\_000)}, takes its cumulative
sum, and compares \code{np.std} of the oscillator against $s$ evaluated
directly.

\bibliographystyle{plainnat}
\bibliography{references}

\end{document}
